\documentclass[11pt]{article}
\usepackage[margin=1in]{geometry}
\usepackage{amsmath,amssymb,booktabs,longtable,hyperref,array}
\usepackage[T1]{fontenc}
\usepackage{lmodern}
\usepackage{microtype}
\hypersetup{colorlinks=true,linkcolor=black,urlcolor=blue,citecolor=black}
\setlength{\parindent}{0pt}
\setlength{\parskip}{5pt}

\title{A Ternary $(4k+1(2))/3$ Collatz-Type Map with Two Attracting Cycles:\\
Theoretical Analysis and Exhaustive Computational Verification up to $10^9$}
\author{Farhad Banazadeh\\
Independent Researcher\\
\href{mailto:fn.bana@hotmail.com}{fn.bana@hotmail.com}\\
\href{https://orcid.org/0009-0004-7023-0298}{ORCID: 0009-0004-7023-0298}}
\date{3 September 2026}

\begin{document}
\maketitle

\begin{abstract}
This paper studies the ternary $(4k+1(2))/3$ Collatz-type map, based on division by $3$ and a complete $3$-adic compression step. The starting values are positive integers not divisible by $3$. For $n=3p+1$ the pre-reduced value is $4p+2$, while for $n=3p+2$ it is $4p+3$. After this step, all factors of $3$ are removed. In the compressed convention used in the computation, the state $3$ is represented by $1$. Two cycles are observed: $C_1=(1,2)$ and $C_7=(7,10,14,19,26,35,47)$. A $3$-adic heuristic based on Haar measure gives the distribution $\Pr(v_3=k)=2/3^{k+1}$ and $\mathbb{E}[v_3]=1/2$. This gives the average logarithmic drift
$\delta=\ln(4/3)-(1/2)\ln 3\approx-0.2616$, corresponding to an average multiplicative factor of about $0.7698$, or about a $23.02\%$ decrease per step on the logarithmic model. This is a heuristic, not a proof of global convergence.

An exhaustive computation of every admissible starting value $1\le n\le10^9$ was performed. There are exactly $666,666,667$ admissible starts. Every one was classified into one of the two observed cycles, with zero unknown cases. The $C_1$ basin contains $21,785,111$ starts ($3.2677666484\%$), and the $C_7$ basin contains $644,881,556$ starts ($96.7322333516\%$). The computation used $42,759,887,447$ steps in total, with mean $64.1398311384$ steps per admissible start. The largest stopping length was $385$, starting from $696,171,200$. The largest value reached by any tested trajectory was $134,701,251,885,711,310$, obtained from $920,435,228$ at step $115$. A maximum of $19$ factors of $3$ was removed in a single step, and the total number of removed factors was $21,683,837,513$. These results lead to a strong computational conjecture that every positive integer eventually enters one of the two cycles and that no other periodic or divergent orbit exists.
\end{abstract}

\section{Introduction}
The classical Collatz problem asks about the long-term behavior of a very simple integer map. Many related maps replace the division by $2$ with other divisors, use several residue classes, or compress an orbit by removing selected prime factors. The present work studies one such ternary system, built around division by $3$ and complete removal of factors of $3$ after each pre-reduced step.

The main difference from the earlier map studied by the author is important. The earlier map used the two branches $(4k-1)/3$ and $(4k+1)/3$ together with the direct $k/3$ branch. The present map instead uses $4k+2$ for the residue class $1\pmod 3$ and $4k+1$ for the residue class $2\pmod 3$, followed by complete $3$-adic compression. The present system therefore has a different orbit structure and two observed terminal cycles.

The paper has three goals. First, it gives a precise definition of the map and explains why every step is integral. Second, it develops the $3$-adic and average-logarithmic analysis supplied in the theoretical note. Third, it records the large computational experiment through $10^9$, including basin sizes, step counts, peak values, factor-of-$3$ statistics, and record cases.

The computational evidence is deliberately separated from the conjecture. A finite exhaustive computation can establish the behavior of every tested starting value, but it cannot by itself establish the same statement for all positive integers.

\section{Definition of the map}
Let
\[
D=\{n\in\mathbb N:3\nmid n\}.
\]
Every $n\in D$ belongs to exactly one of the two residue classes
\[
n=3p+1\quad\text{or}\quad n=3p+2,
\qquad p\ge0.
\]
Define the pre-reduction map $G:D\to\mathbb N$ by
\[
G(n)=\left\lfloor\frac{4n+2}{3}\right\rfloor.
\]
Equivalently,
\[
G(3p+1)=4p+2,
\qquad
G(3p+2)=4p+3.
\]
The reason for the floor is only to write the two residue cases in one expression; the two displayed formulas show that $G(n)$ is always an integer.

For a positive integer $m$, let
\[
v_3(m)=\max\{r\ge0:3^r\mid m\}
\]
be the exponent of $3$ in $m$. The compressed map is
\[
T(n)=\frac{G(n)}{3^{v_3(G(n))}}.
\tag{1}
\]
Thus $T(n)$ is not divisible by $3$ after the compression.

The computational convention used in the exhaustive program identifies the eliminated state $3$ with the state $1$. In particular,
\[
1\longrightarrow2\longrightarrow1.
\]
This convention is part of the definition of the compressed state space used for the reported computation.

The sequence generated from an initial value $x_0\in D$ is
\[
x_{j+1}=T(x_j),\qquad j\ge0.
\]

\section{How the rule works: numerical examples}
The map is easiest to understand through direct examples.

For $n=1=3(0)+1$,
\[
G(1)=4(0)+2=2,
\qquad T(1)=2.
\]
For $n=2=3(0)+2$,
\[
G(2)=4(0)+3=3=3^1,
\qquad T(2)=1.
\]
Hence $1\to2\to1$.

For $n=4=3(1)+1$,
\[
G(4)=6,
\qquad T(4)=2,
\]
so $4\to2\to1\to2\to1\cdots$.

For $n=5=3(1)+2$,
\[
G(5)=7,
\qquad T(5)=7.
\]
Therefore $5$ enters the seven-cycle immediately.

A central example is
\[
7\to10\to14\to19\to26\to35\to47\to7.
\tag{2}
\]
Every step from $7$ through $35$ has no factor of $3$ to remove. At $47$,
\[
G(47)=4(15)+3=63=3^2\cdot7,
\]
and therefore
\[
T(47)=63/9=7.
\]
This closes the cycle.

Another example showing compression is
\[
2\to3\to1\to2\to1\cdots,
\]
where the state $3$ is normalized to $1$. The compression is important: it makes the dynamics live on the nonmultiples of $3$ and records the complete $3$-adic removal in a single step.

\section{The two observed cycles}
The first observed cycle is
\[
C_1=(1,2),
\]
with length $2$. Its exact transitions are
\[
1\to2,
\qquad 2\to1.
\]

The second observed cycle is
\[
C_7=(7,10,14,19,26,35,47),
\]
with length $7$.

For the seven-cycle, the first six transitions are
\[
7\to10\to14\to19\to26\to35\to47,
\]
and the final transition is
\[
47\to63\to7
\]
in the uncompressed description, or simply $47\to7$ in the compressed map. Thus the final step removes exactly two factors of $3$.

The theoretical note gives multiplicative cycle-scale estimates
\[
\lambda_1=\frac{(4/3)^2}{3}=\frac{16}{27}\approx0.5926<1,
\]
and
\[
\lambda_2=\frac{(4/3)^7}{3^2}
=\frac{16384}{19683}
\approx0.8324<1.
\]
These values are useful heuristic stability indicators for the asymptotic multiplicative model. They should not be confused with ordinary derivatives of the integer map, because the actual map contains the discrete valuation operation.

For comparison, the exact product of the actual ratios around $C_1$ is
\[
\frac21\cdot\frac13=\frac23<1,
\]
while the exact product around $C_7$ telescopes to $1$ because the cycle returns to its starting value. The values $16/27$ and $16384/19683$ are therefore best understood as the theoretical-note's asymptotic scale estimates rather than exact cycle multipliers.

\section{Stopping time and peak}
For an initial value $x_0=n$, write its orbit as
\[
x_0,x_1,x_2,\ldots,
\qquad x_{j+1}=T(x_j).
\]
For the two-cycle system, a natural stopping time is the first index at which the orbit becomes smaller than its initial value:
\[
\sigma(n)=\min\{j\ge1:x_j<n\},
\tag{3}
\]
when such an index exists. This is the stopping-time definition used in the theoretical analysis.

For an orbit segment up to a stopping time, define the peak by
\[
P_\sigma(n)=\max_{0\le j\le\sigma(n)}x_j.
\tag{4}
\]
For the large $10^9$ computation, an additional global trajectory peak was tracked independently of this stopping-time definition. It is the maximum value reached before the trajectory was classified into one of the observed cycles.

The distinction matters. An orbit may have a relatively short stopping time but still make a very large temporary excursion. This is exactly what happens in the largest-peak record reported below.

\section{Methodology}
Two complementary approaches were used.

\subsection{Theoretical $3$-adic analysis}
The first approach studies the distribution of the number of factors of $3$ removed after the pre-reduction step. The analysis uses the Haar-measure heuristic on the ring of $3$-adic integers $\mathbb Z_3$ and assumes that large orbit values sample residue classes in an approximately uniform way.

\subsection{Exhaustive computation}
The second approach is direct computation. The final C audit checks every integer
\[
1\le n\le10^9,
\qquad 3\nmid n.
\]
The number of admissible starts is
\[
10^9-\left\lfloor\frac{10^9}{3}\right\rfloor
=666,666,667.
\]
The final C program uses 64-bit integer arithmetic, checks the two observed cycles explicitly, counts the number of removed factors of $3$, and keeps an explicit unknown counter. The reported $10^9$ result is the C audit result supplied with this work.

The code was run as an exhaustive audit. The reported runtime was $527.834688$ seconds, about $8$ minutes and $48$ seconds, on the author's 2012 MacBook Pro.

\section{Theoretical analysis: average logarithmic contraction}
For large $n$, the pre-reduced value has the approximate scale
\[
G(n)\approx\frac43n.
\]
If
\[
v=v_3(G(n)),
\]
then after complete removal of factors of $3$,
\[
T(n)=\frac{G(n)}{3^v}
\approx n\frac{4}{3^{v+1}}.
\tag{5}
\]
Therefore the logarithmic change is approximately
\[
\ln\frac{T(n)}n
\approx \ln\frac43-v\ln3.
\tag{6}
\]

\subsection{Ergodic and Haar-measure viewpoint}
The theoretical note treats the large-scale orbit as sampling residue classes in an approximately uniform way. In the $3$-adic setting, the natural reference measure is Haar measure on $\mathbb Z_3$. This does not mean that a deterministic orbit is literally random; rather, it gives a useful ergodic model for the frequency of residue classes and therefore for the frequency of powers of $3$ removed by the compression.

The key quantity is the $3$-adic valuation of the pre-reduced value after the compulsory division by $3$. Under the uniform-residue or Haar-measure heuristic in $\mathbb Z_3$,
\[
\Pr(v=0)=1-\frac13=\frac23,
\]
\[
\Pr(v=1)=\frac13-\frac19=\frac29,
\]
and, in general,
\[
\Pr(v=k)=\frac{2}{3^{k+1}},
\qquad k\ge0.
\tag{7}
\]
Indeed, $v=k$ means divisibility by $3^k$ but not by $3^{k+1}$. The probability of divisibility by $3^k$ is $3^{-k}$, and the difference between successive probabilities gives (7).

\subsection{Expected number of removed factors}
Using (7),
\[
\mathbb E[v]
=\sum_{k=0}^{\infty}k\frac{2}{3^{k+1}}
=\frac{2}{9}\sum_{k=1}^{\infty}k\left(\frac13\right)^{k-1}.
\]
Since
\[
\sum_{k=1}^{\infty}kr^{k-1}=\frac{1}{(1-r)^2},
\qquad |r|<1,
\]
we obtain
\[
\mathbb E[v]
=\frac{2}{9}\frac{1}{(1-1/3)^2}
=\frac12.
\tag{8}
\]
Thus the heuristic predicts one half of a factor-of-$3$ exponent on average per step.

\subsection{Logarithmic drift}
Taking the expectation in (6),
\[
\delta
=\mathbb E\left[\ln\frac{T(n)}n\right]
\approx \ln\frac43-\mathbb E[v]\ln3.
\]
Using (8),
\[
\boxed{\delta=\ln\frac43-\frac12\ln3
=\ln\left(\frac4{3\sqrt3}\right)}.
\tag{9}
\]
Numerically,
\[
\delta\approx-0.2616<0.
\]
Therefore the corresponding geometric mean factor is
\[
e^\delta=\frac4{3\sqrt3}\approx0.7698.
\tag{10}
\]
In this model, the orbit size decreases by about
\[
1-0.7698\approx0.2302,
\]
or $23.02\%$, per step on the logarithmic average scale.

This is the central theoretical reason for expecting contraction. The pre-reduction step multiplies the size by roughly $4/3$, but the average removal of $3$-adic factors contributes an additional factor of roughly $3^{-1/2}$. The combined average factor is below $1$.

\subsection{Average multiplicative factor}
Equation (10) gives more than a negative logarithmic drift. It gives the geometric mean scale factor
\[
\rho=\exp(\delta)=\frac{4}{3\sqrt3}\approx0.7698.
\tag{14}
\]
Thus the theoretical model predicts
\[
\frac{T(n)}{n}\quad\text{to have geometric mean approximately }0.7698
\]
for large states. Equivalently, the logarithmic size is expected to decrease by about $0.2616$ per step, while the corresponding geometric size decreases by about $23.02\%$ per step. This is the main contraction mechanism proposed in the theoretical analysis.

\subsection{What the negative drift does and does not prove}
The negative value in (9) is a strong heuristic indication of downward behavior, but it is not a proof of convergence of every orbit. The valuation distribution was obtained from an asymptotic uniformity assumption, while an actual orbit is deterministic and its successive residue classes are not independent random samples.

In particular, the statement that divergence has probability zero belongs to the probabilistic heuristic. It should not be stated as a theorem about every integer orbit. A rigorous proof would require a deterministic argument controlling all possible long-term residue patterns, or a different global argument.

\section{Algebraic view of the seven-cycle}
The seven-cycle is especially simple because the first six compressed steps have no factor of $3$ removed. Its values satisfy
\[
7\mapsto10\mapsto14\mapsto19\mapsto26\mapsto35\mapsto47.
\]
The final pre-reduced value is
\[
G(47)=63=3^2\cdot7.
\]
Therefore
\[
T(47)=\frac{63}{3^2}=7.
\]
The cycle closes only because the final step removes two powers of $3$ at once.

This gives a useful picture of the dynamics. The increasing branches can produce a long upward run, but a sufficiently divisible pre-reduced value can produce a strong downward jump. The seven-cycle is a compact example of this balance: six consecutive increases are followed by one compressed return.

\section{Computational verification up to $10^9$}
The final exhaustive run tested every admissible starting value through $10^9$. The result is complete for the finite domain: there were no unresolved starts.

\begin{table}[h]
\centering
\caption{Complete computational result for $1\le n\le10^9$, $3\nmid n$.}
\begin{tabular}{lr}
\toprule
Quantity & Value\\
\midrule
Upper bound $N$ & 1,000,000,000\\
Admissible starts & 666,666,667\\
Cycle $C_1=(1,2)$ count & 21,785,111\\
Cycle $C_1$ percentage & 3.267766648366117\%\\
Cycle $C_7$ count & 644,881,556\\
Cycle $C_7$ percentage & 96.732233351633880\%\\
Unknown & 0\\
Total steps & 42,759,887,447\\
Mean steps & 64.139831138430083\\
Maximum steps & 385\\
Start attaining maximum steps & 696,171,200\\
Global peak & 134,701,251,885,711,310\\
Start attaining global peak & 920,435,228\\
Step of global peak & 115\\
Total removed factors of $3$ & 21,683,837,513\\
Maximum removed in one step & 19\\
Runtime & 527.834688 s\\
\bottomrule
\end{tabular}
\end{table}

As a basic consistency check,
\[
21,785,111+644,881,556=666,666,667,
\]
so the two cycle counts account for every admissible starting value. Since the unknown count is zero, the computation gives a complete classification of the tested domain.

The basin proportions are notably unequal. About $3.27\%$ of the admissible starts enter $C_1$, while about $96.73\%$ enter $C_7$. Thus, in the finite experiment, the seven-cycle is by far the dominant attractor.

\section{A one-million range stopping-time study}
The theoretical note also contains a detailed stopping-time study over the first million integers not divisible by $3$. There are $666,667$ admissible starts in that range. The reported maximum stopping time is
\[
143
\]
at
\[
n=497,116,
\]
with a peak of
\[
410,683,162
\]
within the stopping-time segment.

The ten leading stopping-time records reported in that study are shown below.

\begin{table}[h]
\centering
\caption{Top stopping-time records in the $10^6$ study.}
\begin{tabular}{rrrr}
\toprule
Rank & Start $n$ & $\sigma(n)$ & Peak\\
\midrule
1 & 497,116 & 143 & 410,683,162\\
2 & 662,822 & 142 & 410,683,162\\
3 & 883,763 & 141 & 410,683,162\\
4 & 749,897 & 129 & 2,687,932,606\\
5 & 629,449 & 127 & 263,207,404,847\\
6 & 839,266 & 126 & 263,207,404,847\\
7 & 832,280 & 121 & 2,687,932,606\\
8 & 981,244 & 120 & 40,062,776,062\\
9 & 360,637 & 119 & 517,652,138\\
10 & 999,863 & 118 & 2,687,932,606\\
\bottomrule
\end{tabular}
\end{table}

The value $497,116$ is the stopping-time champion in the one-million study. It survives for 143 steps before first becoming smaller than its starting value, while its orbit reaches about $410$ million in the measured segment. This is a useful example of why stopping time and peak size are different quantities.

The one-million experiment reported approximately $3.31\%$ of starts entering $C_1$ and $96.69\%$ entering $C_7$. The larger $10^9$ audit gives the more precise basin proportions reported in Table 1. These percentages are therefore reported separately rather than mixed together.

\section{The $10^9$ record for maximum stopping length}
The final $10^9$ audit found a maximum stopping length of
\[
385
\]
steps, attained by
\[
n=696,171,200.
\]
This record is different from the one-million stopping-time champion. The one-million study uses the stopping-time definition in (3), while the $10^9$ audit also tracks the number of steps needed to reach a known terminal cycle. The two statistics answer different questions and should not be conflated.

The existence of a record of 385 steps in the billion-scale scan shows that some starting values can resist rapid descent for hundreds of iterations even though all tested trajectories eventually enter one of the two cycles.

\section{The $10^9$ record for largest peak}
The largest trajectory value found in the final audit is
\[
134,701,251,885,711,310,
\]
obtained from the starting value
\[
920,435,228
\]
at step
\[
115.
\]
In scientific notation this is approximately
\[
1.3470125188571131\times10^{17}.
\]

This record is important because it demonstrates that a finite orbit can grow enormously before returning to an attracting cycle. The global peak is therefore much larger than the starting value, even though the starting value is below $10^9$.

The result is also a warning against arguments based only on the first few steps. Temporary growth can be very large. A global proof must control such excursions for arbitrary starting values.

\section{Factor-of-$3$ statistics}
The computation counted the exact number of factors of $3$ removed during all compressed steps. Across all admissible starts up to $10^9$, the total was
\[
21,683,837,513.
\]
The largest number removed in a single step was
\[
19.
\]
Thus at least one tested transition contained a factor $3^{19}$ in the pre-reduced value before compression.

These statistics give direct computational information about the mechanism behind the theoretical contraction model. The average-log calculation treats $v_3(G(n))$ statistically; the exhaustive audit measures how often and how strongly this mechanism actually occurs across the tested trajectories.

\section{Why the two cycles are natural candidates}
The first cycle is created by a short interaction with the factor $3$:
\[
1\to2\to3\to1
\]
in the uncompressed description, which becomes $1\to2\to1$ after normalization.

The second cycle contains a longer increasing chain:
\[
7\to10\to14\to19\to26\to35\to47,
\]
followed by the strong compression
\[
47\to63\to7.
\]
The cycle therefore shows exactly the type of competition suggested by the $3$-adic heuristic: several steps with approximate factor $4/3$ can be balanced by a later removal of powers of $3$.

The $10^9$ basin counts show that the second cycle captures the great majority of tested starts. This does not mean that $C_7$ is globally dominant in a mathematical sense; it means only that it is dominant in the finite computational domain tested here.

\section{Computational conjecture}
The numerical evidence leads to the following conjecture.

\textbf{Conjecture (Two-Cycle Convergence Conjecture).} Every positive integer, under the compressed ternary $3$-adic map defined by (1), eventually enters one of the two cycles
\[
C_1=(1,2)
\]
and
\[
C_7=(7,10,14,19,26,35,47).
\]
In particular, no other periodic cycle and no divergent orbit to infinity exists in the positive integers.

The exhaustive computation through $10^9$ provides strong finite evidence for this conjecture. It does not prove the conjecture for all positive integers.

\section{What is still needed for a proof?}
A proof of the conjecture would have to address at least three main issues.

First, the negative drift in (9) depends on a Haar-measure or uniform-residue heuristic. A rigorous proof would need to control the deterministic frequency of $3$-adic valuations along every orbit, or avoid the probabilistic model entirely.

Second, a proof must rule out every other periodic cycle. The computation found only $C_1$ and $C_7$ in the tested range, but a cycle whose entries lie above the tested range is not excluded by finite computation.

Third, a proof must rule out an orbit that grows forever. The global peak of more than $10^{17}$ from a start below $10^9$ shows that very large temporary growth is possible. Therefore a proof cannot rely only on small numerical behavior.

These are genuine open parts of the problem as studied here. The paper therefore treats the theoretical contraction as a heuristic and the $10^9$ result as a finite theorem about the tested set, while treating the all-integers statement as a conjecture.

\section{Relation to the author's earlier ternary map}
The present map is related to the author's earlier paper, but it is not the same map. The earlier work studied the ternary $(4k\pm1)/3$ map with a fixed point at $1$. The present map uses the pre-reduction rule
\[
G(3p+1)=4p+2,
\qquad
G(3p+2)=4p+3,
\]
followed by complete removal of factors of $3$, and it has the two observed cycles $C_1$ and $C_7$.

The earlier article is therefore cited as related background, not as a source of the present computational result. The two studies should be kept as separate records because their maps, terminal structures, and numerical behavior are different.

The author's previous article is available through Zenodo at
\href{https://doi.org/10.5281/ZENODO.22195651}{doi:10.5281/ZENODO.22195651},
and the computational repository is
\href{https://github.com/FarhadBanazadeh/sequence-computations}{github.com/FarhadBanazadeh/sequence-computations}.

\section{Reproducibility}
The final computational package should contain the LaTeX source, the PDF, the exhaustive C++ source, the numerical result files, and a short README. The C program records the following items: admissible-start count, basin counts and percentages, unknown count, total steps, mean steps, maximum steps and its starting value, global peak and its starting value and step, total removed factors of $3$, maximum number of factors removed in one step, and runtime.

The main computation was performed with the supplied C program. A standard compilation command is
\begin{verbatim}
cc -O3 -std=c11 ternary_collatz.c -o ternary_collatz
./ternary_collatz
\end{verbatim}
The reported run time was $527.834688$ seconds on a 2012 MacBook Pro.

The finite computation is deterministic. The final C audit follows each admissible start until one of the two known cycles is reached or the safety step limit is reached, while retaining aggregate statistics for the full run.

\section{Conclusion}
This paper introduced a new ternary $3$-adic Collatz-type map in which the pre-reduced value is approximately $4n/3$ and all factors of $3$ are then removed. The resulting compressed dynamics have two observed cycles, $C_1=(1,2)$ and $C_7=(7,10,14,19,26,35,47)$.

The theoretical analysis gives a simple explanation for why contraction is plausible. Under the Haar-measure heuristic,
\[
\Pr(v_3=k)=\frac{2}{3^{k+1}},
\qquad
\mathbb E[v_3]=\frac12,
\]
and hence
\[
\delta=\ln\frac43-\frac12\ln3\approx-0.2616.
\]
The corresponding average multiplicative factor is about $0.7698$, suggesting a $23.02\%$ decrease per step on the logarithmic scale. This is useful evidence, but it is not a proof.

The computational evidence is much stronger in the finite domain. Every one of the $666,666,667$ admissible starting values up to $10^9$ was classified into one of the two observed cycles, with zero unknown cases. The $C_1$ basin contains $21,785,111$ starts, while $C_7$ contains $644,881,556$. The total number of recorded steps is $42,759,887,447$, the mean is $64.1398311384$, the maximum stopping length reported by the final audit is $385$, and the largest trajectory peak is $134,701,251,885,711,310$.

The results strongly support the Two-Cycle Convergence Conjecture: every positive integer eventually enters $C_1$ or $C_7$, with no other cycle and no divergent orbit. The conjecture remains open. The purpose of the present paper is to provide a clear mathematical definition, a complete theoretical analysis of the $3$-adic mechanism, a large finite computation, and enough source and data material for independent reproduction.

\section*{Data and code availability}
The supplementary archive associated with this paper contains the exhaustive verification source, the reported numerical results, and documentation of the computation. The paper source is supplied separately so that the PDF and LaTeX source can be deposited independently on Zenodo and GitHub.

\begin{thebibliography}{9}
\bibitem{lagarias1985}
J. C. Lagarias, ``The 3x+1 problem and its generalizations,'' \emph{The American Mathematical Monthly}, 92(1), 3--23, 1985. DOI: \href{https://doi.org/10.2307/2322189}{10.2307/2322189}.

\bibitem{banazadeh2026}
Farhad Banazadeh, ``A Ternary $(4k\pm1)/3$ Collatz-Type Map: Computational Verification up to $10^9$,'' 2026. Zenodo DOI: \href{https://doi.org/10.5281/ZENODO.22195651}{10.5281/ZENODO.22195651}.

\bibitem{github}
Farhad Banazadeh, ``Sequence Computations,'' GitHub repository, 2026. \href{https://github.com/FarhadBanazadeh/sequence-computations}{github.com/FarhadBanazadeh/sequence-computations}.
\end{thebibliography}

\end{document}
